\documentclass[12pt,a4paper]{article}
\usepackage{fontspec}
\usepackage[russian]{babel}
\usepackage{amsmath,amssymb,amsfonts,amsthm}
\usepackage{geometry}
\usepackage{enumitem}
\usepackage{microtype}

\geometry{margin=2cm}
\defaultfontfeatures{Ligatures=TeX}
\setmainfont{Times New Roman}
\setlength{\parindent}{0pt}
\setlength{\parskip}{0.45em}
\setlist[itemize,enumerate]{itemsep=0.55em, topsep=0.7em, parsep=0.3em}
\setlength{\abovedisplayskip}{0.8em}
\setlength{\belowdisplayskip}{0.8em}

\begin{document}

\section*{Задание 2. Кривые}

\textit{Ниже используются стандартные параметризации: окружность радиуса $R$, винтовая линия $\gamma(t) = (a\cos t, a\sin t, bt)$ и график $\gamma(t) = (t, f(t), 0)$.}

\subsection*{Формулы, которые используются}

\begin{enumerate}
    \item \textbf{Вектор скорости.} Если
    \[
    \gamma(t) = (x(t), y(t), z(t)),
    \]
    то
    \[
    \gamma'(t) = (x'(t), y'(t), z'(t)), \qquad |\gamma'(t)| = \sqrt{x'(t)^2 + y'(t)^2 + z'(t)^2}.
    \]

    \item \textbf{Касательная прямая.} В точке $t_0$:
    \[
    (X,Y,Z) = \gamma(t_0) + \lambda\,\gamma'(t_0).
    \]

    \item \textbf{Нормальная плоскость.} В точке $t_0$:
    \[
    \gamma'(t_0)\cdot\bigl((X,Y,Z)-\gamma(t_0)\bigr)=0.
    \]

    \item \textbf{Длина дуги.}
    \[
    L(\gamma) = \int_{\alpha}^{\beta} |\gamma'(t)|\,dt.
    \]

    \item \textbf{Естественная параметризация.}
    \[
    s(t) = \int_{t_0}^{t} |\gamma'(\tau)|\,d\tau, \qquad \left\| \frac{d\tilde{\gamma}}{ds} \right\| = 1.
    \]

    \item \textbf{Базис Френе.}
    \[
    \vec{\tau} = \frac{d\tilde{\gamma}}{ds}, \qquad
    \vec{n} = \frac{1}{k}\frac{d\vec{\tau}}{ds}, \qquad
    \vec{b} = \vec{\tau} \times \vec{n}.
    \]

    \item \textbf{Кривизна и кручение.}
    \[
    k = \left\| \frac{d\vec{\tau}}{ds} \right\|, \qquad
    \rho = \frac{1}{k}, \qquad
    \varkappa = -\left( \frac{d\vec{b}}{ds}, \vec{n} \right).
    \]

    \item \textbf{График функции $y=f(x)$.}
    Для $\gamma(t) = (t, f(t), 0)$:
    \[
    k = \frac{|f''(t)|}{(1+(f'(t))^2)^{3/2}}.
    \]
\end{enumerate}

\subsection*{1. Окружность}

Пусть
\[
\gamma(t) = (R\cos t, R\sin t, 0), \qquad R > 0.
\]

\textbf{1) Вектор скорости. Касательная прямая и нормальная плоскость}

\[
\gamma'(t) = (-R\sin t, R\cos t, 0), \qquad |\gamma'(t)| = R.
\]

Касательная прямая в точке $t_0$:
\[
(x,y,z) = (R\cos t_0, R\sin t_0, 0) + \lambda(-R\sin t_0, R\cos t_0, 0).
\]

Нормальная плоскость в точке $t_0$:
\[
-\sin t_0\,(x - R\cos t_0) + \cos t_0\,(y - R\sin t_0) = 0,
\]
то есть
\[
-x\sin t_0 + y\cos t_0 = 0.
\]

\textbf{2) Естественная параметризация и длина дуги}

\[
s(t) = \int_0^t |\gamma'(\tau)|\,d\tau = Rt, \qquad t = \frac{s}{R}.
\]

Естественная параметризация:
\[
\tilde{\gamma}(s) = \left(R\cos \frac{s}{R}, R\sin \frac{s}{R}, 0\right).
\]

Длина дуги на отрезке $[t_1,t_2]$:
\[
L = \int_{t_1}^{t_2} |\gamma'(t)|\,dt = R|t_2 - t_1|.
\]

Если берется вся окружность, то
\[
L = 2\pi R.
\]

\textbf{3) Базис Френе, кривизна, радиус кривизны, кручение}

\[
\vec{\tau} = \left(-\sin \frac{s}{R}, \cos \frac{s}{R}, 0\right), \qquad
\vec{n} = \left(-\cos \frac{s}{R}, -\sin \frac{s}{R}, 0\right), \qquad
\vec{b} = (0,0,1).
\]

\[
k = \left\| \frac{d\vec{\tau}}{ds} \right\| = \frac{1}{R}, \qquad
\rho = \frac{1}{k} = R, \qquad
\varkappa = 0.
\]

\subsection*{2. Винтовая линия}

Пусть
\[
\gamma(t) = (a\cos t, a\sin t, bt), \qquad a > 0.
\]

\textbf{1) Вектор скорости. Касательная прямая и нормальная плоскость}

\[
\gamma'(t) = (-a\sin t, a\cos t, b), \qquad |\gamma'(t)| = \sqrt{a^2+b^2}.
\]

Касательная прямая в точке $t_0$:
\[
(x,y,z) = (a\cos t_0, a\sin t_0, bt_0) + \lambda(-a\sin t_0, a\cos t_0, b).
\]

Нормальная плоскость в точке $t_0$:
\[
-a\sin t_0\,(x-a\cos t_0) + a\cos t_0\,(y-a\sin t_0) + b(z-bt_0) = 0.
\]

\textbf{2) Естественная параметризация и длина дуги}

\[
s(t) = \int_0^t |\gamma'(\tau)|\,d\tau = \sqrt{a^2+b^2}\,t,
\qquad
t = \frac{s}{\sqrt{a^2+b^2}}.
\]

Естественная параметризация:
\[
\tilde{\gamma}(s) = \left(
a\cos \frac{s}{\sqrt{a^2+b^2}},
a\sin \frac{s}{\sqrt{a^2+b^2}},
\frac{bs}{\sqrt{a^2+b^2}}
\right).
\]

Длина дуги на отрезке $[t_1,t_2]$:
\[
L = \int_{t_1}^{t_2} |\gamma'(t)|\,dt = \sqrt{a^2+b^2}\,|t_2-t_1|.
\]

\textbf{3) Базис Френе, кривизна, радиус кривизны, кручение}

Положим
\[
\xi = \frac{s}{\sqrt{a^2+b^2}}.
\]
Тогда
\[
\vec{\tau} = \left(
\frac{-a}{\sqrt{a^2+b^2}}\sin \xi,
\frac{a}{\sqrt{a^2+b^2}}\cos \xi,
\frac{b}{\sqrt{a^2+b^2}}
\right),
\]
\[
\vec{n} = (-\cos \xi, -\sin \xi, 0),
\]
\[
\vec{b} = \left(
\frac{b}{\sqrt{a^2+b^2}}\sin \xi,
-\frac{b}{\sqrt{a^2+b^2}}\cos \xi,
\frac{a}{\sqrt{a^2+b^2}}
\right).
\]

\[
k = \left\| \frac{d\vec{\tau}}{ds} \right\| = \frac{a}{a^2+b^2}, \qquad
\rho = \frac{1}{k} = \frac{a^2+b^2}{a}, \qquad
\varkappa = \frac{b}{a^2+b^2}.
\]

\subsection*{3. График функции $y=f(x)$}

Пусть
\[
\gamma(t) = (t, f(t), 0), \qquad f \in C^2.
\]

\textbf{1) Вектор скорости. Касательная прямая и нормальная плоскость}

\[
\gamma'(t) = (1, f'(t), 0), \qquad |\gamma'(t)| = \sqrt{1 + (f'(t))^2}.
\]

Касательная прямая в точке $t_0$:
\[
(x,y,z) = (t_0, f(t_0), 0) + \lambda(1, f'(t_0), 0).
\]

Нормальная плоскость в точке $t_0$:
\[
(x-t_0) + f'(t_0)\bigl(y-f(t_0)\bigr) = 0.
\]

\textbf{2) Естественная параметризация и длина дуги}

\[
s(t) = \int_{t_0}^{t} \sqrt{1 + (f'(\tau))^2}\,d\tau.
\]

Если $t=t(s)$ --- обратная функция, то естественная параметризация имеет вид
\[
\tilde{\gamma}(s) = \bigl(t(s), f(t(s)), 0\bigr).
\]

Длина дуги на отрезке $[t_1,t_2]$:
\[
L = \int_{t_1}^{t_2} \sqrt{1 + (f'(t))^2}\,dt.
\]

\textbf{3) Базис Френе, кривизна, радиус кривизны, кручение}

В точках, где $f''(t_0) \neq 0$, можно взять
\[
\vec{\tau} = \frac{(1, f'(t), 0)}{\sqrt{1+(f'(t))^2}},
\qquad
\vec{n} = \operatorname{sgn}(f''(t)) \frac{(-f'(t), 1, 0)}{\sqrt{1+(f'(t))^2}},
\qquad
\vec{b} = \operatorname{sgn}(f''(t))(0,0,1).
\]

\[
k = \frac{|f''(t)|}{\bigl(1 + (f'(t))^2\bigr)^{3/2}}, \qquad
\rho = \frac{\bigl(1 + (f'(t))^2\bigr)^{3/2}}{|f''(t)|}, \qquad
\varkappa = 0.
\]

Если $f''(t_0)=0$, то в этой точке $\kappa=0$, базис Френе не определён, и кручение: нет ответа.

\end{document}
